\usepackage[papersize={24.4602cm,28.575cm},textwidth=18.9cm,
textheight=23.5cm,topmargin=30mm,botmargin=20.75mm,headsep=0mm,headheight=0mm,
leftmargin=28mm,rightmargin=28mm,footskip=0pt]{zwpagelayout}

%%try following
%The standard book size closest to 15.4 cm \(\times \) 19 cm is Crown Quarto 
%(6.84" \(\times \) 8.5" or 17.1 cm \(\times \) 21.6 cm), 
%though it is slightly larger. In localized school supplies, 
%it directly matches the Small Four-Line or Single-Line Notebook dimensions 
%(usually 15.5 cm \(\times \) 19 cm) commonly used in Indi


%~ \usepackage[
  %~ paperwidth=28.575cm,
  %~ paperheight=24.4602cm,
  %~ top=2cm,
  %~ bottom=2.5cm,
  %~ inner=2.5cm,
  %~ outer=2cm,
  %~ headsep=0.5cm,
  %~ footskip=1cm,
  %~ includehead,
  %~ includefoot,
  %~ showcrop,
%~ ]{geometry}

\usepackage{polyglossia}
\usepackage{xltxtra}
\usepackage{fontspec}
\usepackage{endnotes}
\usepackage{graphicx} 
\usepackage{wrapfig}
\usepackage{setspace}
\usepackage{multicol}
\usepackage{xcolor}
%\usepackage{pagecolor}
\usepackage{calc}
\usepackage{fancyhdr}
\usepackage{enumitem}
\usepackage{stackengine}
\usepackage{changepage} 
%\usepackage{showframe}

%testing for kannada toc
%using regular toc file for english
\usepackage{titletoc}

%\definecolor{ArchivalCream}{HTML}{f9f1f1} % A soft cream for your archive
%\definecolor{ArchivalCream}{HTML}{F5EFE2} % A soft cream for your archive

%color suggested by rajaneesha
%\definecolor{ArchivalCream}{HTML}{F4E7D6} % A soft cream for your archive
%\definecolor{ArchivalCream}{HTML}{EFE5D6} % A soft cream for your archive
%\definecolor{ArchivalCream}{HTML}{F1E1D2} % A soft cream for your archive

%\definecolor{ArchivalCream}{HTML}{F4EFE7} % A soft cream for your archive
%\pagecolor{ArchivalCream}

%\definecolor{deepmaroon}{RGB}{193,154,107}
\definecolor{deepmaroon}{RGB}{181,143,95}
\definecolor{markcolor}{RGB}{181,143,95}
\definecolor{bordercolor}{RGB}{215, 75, 15}
%~ \definecolor{plum}{RGB}{153, 101, 21}
%\definecolor{plum}{RGB}{6, 1, 110}
\definecolor{plum}{RGB}{235, 235, 235}

\usepackage{crop}
\usepackage{eso-pic}
\usepackage{tikz}
\usepackage{adjustbox}
\usetikzlibrary{calc}
\usepackage{amssymb}
%\usepackage{MnSymbol}
% Define Border Properties (same as before)
\newlength{\borderthickness}
\setlength{\borderthickness}{0.05mm} % Set the desired thickness of your border line

\newlength{\borderoffset}
\setlength{\borderoffset}{10mm} % Set the distance of the border from the page edge



% Main language
\setmainlanguage{english}

% Other languages
\setotherlanguages{kannada, sanskrit}

% Font settings for each language
\setmainfont[Path=fonts/,
    Scale=1.1,
    ItalicFont=Junicode-Italic.otf,
    %SmallCapsFont=Junicode-Regular.otf,
    BoldFont=Junicode-Bold.otf,
    BoldItalicFont=Junicode-BoldItalic.otf,
    Numbers=OldStyle]{Junicode-Regular.otf} 

\newfontfamily\englishfont[Path=fonts/,
    Scale=1.1,
    ItalicFont=Junicode-Italic.otf,
    %SmallCapsFont=Junicode-Regular.otf,
    BoldFont=Junicode-Bold.otf,
    BoldItalicFont=Junicode-BoldItalic.otf,
    Numbers=OldStyle]{Junicode-Regular.otf}

\newfontfamily\montsefont[Path=fonts/,
    Scale=1.1,
    ItalicFont=Montserrat-BlackItalic.otf,
    BoldFont=Montserrat-Bold.otf,
    BoldItalicFont=Montserrat-BoldItalic.otf
    ]{Montserrat-Black.otf}

\newfontfamily\kannadafont[
	Script=Kannada,
	Scale=1.2,
	BoldFont=SHREE-KAN-OTF-0850-Bold,
	ItalicFont=SHREE-KAN-OTF-0850-Italic,
	BoldItalicFont=SHREE-KAN-OTF-0850-Bold-Italic,
	HyphenChar="200C
]{SHREE-KAN-OTF-0850}

\newfontfamily\sanskritfont 
[   Script=Devanagari,
	Scale=1.4,
    Mapping=RomDev]{Adishila}

\newfontfamily\tmpfont 
[   Script=Devanagari,
	Scale=1.4
    ]{Adishila}

\setlength{\columnsep}{1cm}

\makeatletter
\renewcommand\chapter{%\if@openright\cleardoublepage\else\clearpage\fi
					\global\@topnum\z@
                    \@afterindentfalse
                    \secdef\@chapter\@schapter}

\newcommand{\late@tocka}{tocka} % Extension for Kannada ToC
\makeatother

\newenvironment{sans}{\bgroup\selectlanguage{sanskrit}\setstretch{1.4}}{\par\egroup}
\newenvironment{kan}{\bgroup\selectlanguage{kannada}\setstretch{1.4}}{\par\egroup}
\newenvironment{eng}{\bgroup\selectlanguage{english}\setstretch{1.4}}{\par\egroup}

\def\sansinline#1{{\setstretch{1.4}\textsanskrit{#1}}}
\def\kaninline#1{{\setstretch{1.4}\textkannada{#1}}}
\def\enginline#1{{\setstretch{1.4}\textenglish{#1}}}
\newcommand{\kantoc}[1]{{\kannadafont #1}}

\long\def\mainhead#1#2{\ifodd\value{page}%
\preparehead{#1}%
\else
\null\clearpage
\preparehead{#1}
\fi
\addcontentsline{toc}{chapter}{#2}
\newpage
}

\long\def\preparehead#1{%
\null\vfill
\begin{eng}
\begin{center}
{\montsefont\fontsize{16pt}{19.2pt}\fontseries{bx}\selectfont \color{black}{\uppercase{#1}}}	
\end{center}
\end{eng}
\vfill
}

\def\booktitle#1{{\fontsize{36pt}{43.2pt}\fontseries{bx}\selectfont \color{deepmaroon}{#1}}}
\def\titleinline#1{\centerline{\LARGE #1}}

\long\def\prtheading#1{
  \iflanguage{kannada}{%
    % If current language is Kannada
    {\LARGE\bfseries{#1}}%
  }{%
    % If English (default)
    {\montsefont\LARGE\bfseries{#1}}%
  }%
}

\newcommand{\Section}[2][]{%
  \iflanguage{kannada}{%
    % If current language is Kannada
    {\Large\bfseries\color{deepmaroon}{#2}}%
    \addcontentsline{tocka}{section}{\kantoc{\if\relax\detokenize{#1}\relax #2\else #1\fi}}%
  }{%
    % If English (default)
    {\montsefont\large\bfseries\color{deepmaroon}{#2}}%
    \addcontentsline{toc}{section}{\if\relax\detokenize{#1}\relax #2\else #1\fi}%
  }%
}
\newcommand{\Subsection}[2][]{%
  {\large\bfseries #2}%
  \iflanguage{kannada}{%
    % If current language is Kannada, send to Kannada ToC
    \addcontentsline{tocka}{subsection}{\kantoc{\if\relax\detokenize{#1}\relax #2\else #1\fi}}%
  }{%
    % If English (default), send to standard ToC
    \addcontentsline{toc}{subsection}{\if\relax\detokenize{#1}\relax #2\else #1\fi}%
  }%
}
\newcommand{\SSubsection}[2][]{%
  #2%
  \iflanguage{kannada}{%
    % If current language is Kannada, send to Kannada ToC
    \addcontentsline{tocka}{subsubsection}{\kantoc{\if\relax\detokenize{#1}\relax #2\else #1\fi}}%
  }{%
    % If English (default), send to standard ToC
    \addcontentsline{toc}{subsubsection}{\if\relax\detokenize{#1}\relax #2\else #1\fi}%
  }%
}

\def\heading#1{{\large\color{deepmaroon}{\textbf{#1}}}}%
\def\subheading#1{{\textbf{#1}}}%
\def\suggheading#1{{\LARGE\color{deepmaroon}{#1}}}
\def\suggheadingk#1{%
{{\LARGE\iflanguage{kannada}{\color{deepmaroon}{#1}}{#1}}}%
}
\def\headingnoclr#1{%
{\large #1}%
}
\def\genheading#1{\centerline{\large\bfseries #1}}
\def\dummyimage{~\vskip 10mm
\centerline{\fcolorbox{plum}{white}{\phantom{\rule{15cm}{6cm}}}}
\vskip 20mm}
\def\fclrbox#1{\fcolorbox{plum}{white}{#1}}
\def\img{\ignorespaces\centering}
\def\endimg{\vskip 5mm}

\newcounter{anum}
\def\asananum{\enginline{\theanum.}}
\def\incanum{\stepcounter{anum}}

\long\def\insimgleaf#1{
\makebox[\textwidth][c]{%
\fclrbox{\includegraphics[width=\textwidth]{#1}}}
}

\newcommand{\insimg}[2][0.8]{%
\makebox[\textwidth][c]{%
\fclrbox{\includegraphics[width=#1\textwidth]{#2}}%
}}

\long\def\insimgspl#1{
\makebox[\textwidth][c]{%
\fclrbox{\includegraphics[width=\textwidth]{#1}}}
}

%%normalsize heading with color
\def\headingnsize#1{{\color{deepmaroon}{\textbf{#1}}}}

\newenvironment{HalfText}{\begin{minipage}[c][0.49\textheight][c]{\linewidth}}{\end{minipage}}

\newcommand{\HalfPageImage}[2]{%
  \noindent
  \begin{minipage}[c][0.49\textheight][c]{\linewidth}
    \centering
    \insimgleaf{#1}%    
    \Rline{#2}
  \end{minipage}%
}

\makeatletter

\setcounter{secnumdepth}{4}
\setcounter{tocdepth}{4}

\renewcommand\tableofcontents{%
    \vspace*{70\p@}%
    \if@twocolumn
      \@restonecoltrue\onecolumn
    \else
      \@restonecolfalse
    \fi
    {\montsefont\LARGE\bfseries\color{black}{\MakeUppercase\contentsname}}%
    \vskip 40\p@%
    \@starttoc{toc}
    \if@restonecol\twocolumn\fi
    }

\renewcommand{\@dotsep}{10000}
\def\@makeenmark{$^{\theenmark}$}

\renewcommand{\enoteheading}{%
  \section*{\notesname}%
  \markboth{}{} % This clears the header text for this section
}

\renewcommand{\enoteformat}{%
   \noindent\makebox{\theenmark. }% Place number in the margin space
}

\makeatother

\parindent=0pt
\parskip=5pt
%\hfuzz=1pt
\graphicspath{{./images/jpg/}}
\setlength\fboxrule{3pt}
\setlength\fboxsep{0pt}

%~ \renewcommand*\l@chapter[2]{%
  %~ \ifnum \c@tocdepth >\m@ne
    %~ \addpenalty{-\@highpenalty}%
    %~ \vskip 0.7em \@plus\p@
    %~ \setlength\@tempdima{1.5em}%
    %~ \begingroup
      %~ \parindent \z@ \rightskip \@pnumwidth
      %~ \parfillskip -\@pnumwidth
      %~ \leavevmode 
      %~ \advance\leftskip\@tempdima
      %~ \hskip -\leftskip
      %~ \iflanguage{sanskrit}{{\fontsize{15}{18}\selectfont #1}}{%
      %~ \iflanguage{kannada}{{\fontsize{13}{16}\selectfont #1}}{%
      %~ \iflanguage{english}{{\fontsize{11}{14}\selectfont #1}}{}%
      %~ }%
      %~ }\nobreak\hfil \nobreak\hb@xt@\@pnumwidth{\hss #2}\par
      %~ \penalty\@highpenalty
    %~ \endgroup
  %~ \fi}

\long\def\asana#1#2{%
\incanum%
\begin{multicols}{2}
\begin{eng}
\begin{enumerate}[leftmargin=*,labelindent=0pt]
\item[\asananum]
#1
\end{enumerate}
\end{eng}
\begin{kan}
\begin{enumerate}[leftmargin=*,labelindent=0pt]
\item[\asananum]
#2
\end{enumerate}
\end{kan}
\end{multicols}
}

\long\def\asananames#1#2#3{%
\rightline{%
\begin{tabular}{r@{}}
\sansinline{{\small{#1}}}\\ 
\enginline{\small {#2}}\\
\kaninline{\small {#3}}
\end{tabular}}
\smallskip
} 

\long\def\asananamestmp#1#2#3{%
\hbox to\textwidth{\hfil\sansinline{{{\small #1}}}\hspace{0.5cm}
\enginline{{\small #2}}\hspace{0.5cm}
\kaninline{{\small #3}}\hfil}
} 


% Add the Border to Every Page
\AddToShipoutPictureBG{%
\begin{tikzpicture}[overlay, remember picture]
  \node[inner sep=0pt, minimum width=\paperwidth, minimum height=\paperheight] at (current page.center) {
    \includegraphics[width=21.8cm, height=27cm, keepaspectratio]{border2.png}
  };
\end{tikzpicture}%
\begin{tikzpicture}[overlay, remember picture]
    % Placing the page number at coordinate (3, -3) from the top-left corner
    %\node[anchor=north west] at (3, -3) {Page \thepage};

    % Placing the page number 2cm from the bottom-right corner
    %\node[anchor=south east] at ($(current page.south) + (0.5cm, 1cm)$) {\thepage};
    \node[anchor=south, yshift=1cm] at (current page.south) {\enginline{\thepage}};
\end{tikzpicture}
}

\def\engname#1{\textbf{#1}}
\def\engiast#1{\textit{#1}}
\def\insmark#1#2#3{\stackon[#1]{#3}{\hspace{#2}\color{markcolor}{$\lnot$}}}  
\def\krnt{\textit{Kapālakuraṇṭakapaddhati}}
\def\RE{\textit{recto}}
\def\VE{\textit{verso}}
\def\Rline#1{\rightline{\enginline{\small #1}}}
\def\tbf#1{\textbf{#1}}
\def\tit#1{\textit{#1}}

\long\def\insertFullImage#1{%
\ClearShipoutPictureBG % Stops the border/page number for this page
\thispagestyle{empty}

\begin{tikzpicture}[remember picture, overlay]
    % Your "Full Page" content here
    \node[anchor=center, inner sep=0pt] at (current page.center) {
        \includegraphics[width=\paperwidth, height=\paperheight]{sample/#1}
    };
\end{tikzpicture}
}

\long\def\pagesWithBorder{%
\AddToShipoutPictureBG{%
\begin{tikzpicture}[overlay, remember picture]
  \node[inner sep=0pt, minimum width=\paperwidth, minimum height=\paperheight] at (current page.center) {
    \includegraphics[width=21.8cm, height=27cm, keepaspectratio]{border2.png}
  };
\end{tikzpicture}%
\begin{tikzpicture}[overlay, remember picture]
    % Placing the page number at coordinate (3, -3) from the top-left corner
    %\node[anchor=north west] at (3, -3) {Page \thepage};

    % Placing the page number 2cm from the bottom-right corner
    %\node[anchor=south east] at ($(current page.south) + (0.5cm, 1cm)$) {\thepage};
    \node[anchor=south, yshift=1cm] at (current page.south) {\enginline{\thepage}};
\end{tikzpicture}
}
}
  
\pagestyle{fancy}
\fancyhf{}
\renewcommand\headrulewidth{0pt}
\cfoot[]{}
%\overfullrule=15pt

\newcommand{\sethyphenation}[3][]{%
  \sbox0{\begin{otherlanguage}[#1]{#2}
    \hyphenation{#3}\end{otherlanguage}}}

\input dictionary

\brokenpenalty=10000
\widowpenalty=10000
\clubpenalty=10000


\long\def\asanatmp#1{%
\incanum%
\begin{eng}
\begin{enumerate}[leftmargin=*,labelindent=0pt]
\item[\asananum]
#1
\end{enumerate}
\end{eng}
}

\long\def\asanatmpk#1{%
\incanum%
\begin{kan}
\begin{enumerate}[leftmargin=*,labelindent=0pt]
\item[\asananum]
#1
\end{enumerate}
\end{kan}
}
